\subsection{Sobolev 空间的刻画}\label{sec:sobolev-characterizations}

Sobolev 空间的基本定义将其描述为 $L^p(\Omega)$ 的子空间. 当然, 还有许多内容可讨论, 本节将介绍刻画 Sobolev 空间中函数的一些最重要方式.

在本节的大部分内容中, 将把讨论限于 $p=2$ 的情形. 所讨论的许多结果对一般 $p$ 都有类似结论, 对此请读者参阅文献.

\subsubsection{关于区域 $\Omega$ 的一些说明}\label{subsec:sobolev-domain-boundary-hypotheses}

许多有关 Sobolev 空间的问题, 其答案依赖于对 $\Omega$ 的边界正则性的假设.\footnote{在这里对 Sobolev 空间的简短介绍中, 选择避开大多数边界光滑性问题. 关于这一主题的更完整研究, 推荐 Fraenkel 的论文 \cite{Fr}.} 大多数时候将假设边界光滑. 具体地, 作如下定义.

\begin{Definition}
\label{def:c-k-domain-boundary}
若 $\partial\Omega$ 上每个点都有邻域 $N$, 使 $\partial\Omega\cap N$ 为 $C^k$ 曲面, 并且 $\Omega\cap N$ 位于 $\partial\Omega\cap N$ 的 ``一侧'', 则称 $\Omega$ 为 $C^k$ 类, $k\geq1$.
\end{Definition}

若 $\Omega$ 为有界区域, 即连通, 则最后一个假设是多余的; 参见 Remark~\ref{rem:boundary-one-side-jordan}.

在应用中有两类问题涉及非光滑区域:\begin{enumerate}
\item 有角点的区域, 以及
\item 其中 $\Omega$ 事先未知的自由边界问题.
\end{enumerate}事实上, Sobolev 空间的许多结果并不要求光滑边界. 取而代之, 需要假设 ``线段性质'' 和 ``锥性质'' 等各种几何条件 (参见 \cite{Fr}). 这里不讨论这些条件, 但将陈述 Lipschitz 区域的一些结果.

\begin{Definition}
\label{def:lipschitz-domain}
若 $\partial\Omega$ 上每个点都有邻域 $N$, 使在作仿射坐标变换 (平移和旋转) 后, $\partial\Omega\cap N$ 可由方程 $x_m=\phi(x_1,\ldots,x_{m-1})$ 描述, 其中 $\phi$ \textbf{一致 Lipschitz 连续}, 则称 $\Omega$ 为 \textbf{Lipschitz} 区域. 此外, $\Omega\cap N$ 位于 $\partial\Omega\cap N$ 的一侧, 例如 $\Omega\cap N=\{\bm{x}\in N\mid x_m<\phi(x_1,\ldots,x_{m-1})\}$.
\end{Definition}

若 $\Omega$ 无界, 则除了对 $\partial\Omega$ 的光滑性条件以外, 还需要施加表明 $\Omega$ 在无穷远处行为良好的条件. 这里不一般地讨论这些条件, 许多结果将只对 $\partial\Omega$ 有界的情形陈述.

最后, 定义区域 $\Omega$ 的一种刻画, 它作为简洁的技术性假设将十分有用.

\begin{Definition}
\label{def:sobolev-k-extension-property}
若存在有界线性映射 $E:H^k(\Omega)\to H^k(\mathbb{R}^m)$, 使对每个 $u\in H^k(\Omega)$ 有 $Eu|_{\Omega}=u$, 则称 $\Omega$ 具有\textbf{$k$ 阶延拓性质}.
\end{Definition}

反过来, $H^k(\mathbb{R}^m)$ 中每个函数的限制属于 $H^k(\Omega)$, 这当然是显然的. 将在后面的一小节中研究延拓性质; 事实上, 有界 Lipschitz 区域对每个 $k$ 都具有延拓性质.
